\chapter*{使用说明与考情总览}
\addcontentsline{toc}{chapter}{使用说明与考情总览}
\markboth{使用说明与考情总览}{}

\section*{一、考试科目与参考书}

\begin{center}
\begin{tabularx}{\textwidth}{@{}lX@{}}
\toprule
项目 & 内容 \\
\midrule
招生单位 & 西华大学 \\
科目代码与名称 & \textbf{818 工程流体力学} \\
适用专业 & 085800 能源动力（专硕） / 080700 动力工程及工程热物理（学硕） \\
官方参考书 & 赵琴 主编《流体力学与流体机械》，中国水利水电出版社，2016 年 4 月第 1 版，ISBN 978-7-5170-4239-6 \\
初试考查范围 & 该书\textbf{第 1--8 章}（工程流体力学部分）；第 9--14 章（气体动力学、流体机械）为复试范围 \\
\bottomrule
\end{tabularx}
\end{center}

\section*{二、试卷结构（分值分布）}

\begin{center}
\begin{tabularx}{\textwidth}{@{}lccX@{}}
\toprule
题型 & 题量 & 分值 & 主要考法 \\
\midrule
单项选择题 & 10 & 30 分（3 分/题） & 概念辨析、小计算 \\
填空题 & 10 & 30 分（3 分/空） & 名词定义、公式含义、小计算 \\
判断题 & 10 & 20 分（2 分/题） & 概念辨析 \\
简答题 & 3 & 12 分（4 分/题） & 定义、分类、特点（如水头损失分类、均匀流特点） \\
计算题 & 3 & 58 分（15--20 分/题） & 伯努利方程、动量方程、管路水力计算 \\
\bottomrule
\end{tabularx}
\end{center}

\noindent\textbf{2018 年及以前的旧结构}：选择 30 + 填空 30 + 判断 20 + 计算 90（6--7 题），
另有证明题 12--15 分、作图题（仅出现过一次，7 分）。做早期真题时注意题型差异。

\section*{三、出题风格（六条）}

\begin{enumerate}
  \item 重视基础，新题型很少；
  \item \textbf{真题重复率很高}——同一考点会换数字反复考；
  \item 题量适中，计算量适中；
  \item 围绕参考书（赵琴教材）复习即可；
  \item 总体难度中等；
  \item \textbf{每年的真题大部分都源于课后习题的改编}，课后习题必须逐题过关。
\end{enumerate}

\section*{四、逐章重要度（依据 `8-1考情分析.pptx` slide33 与 `9-1考点分析.pptx`）}

\begin{center}
\begin{tabularx}{\textwidth}{@{}clX@{}}
\toprule
章 & 重要度 & 定位（原文口径） \\
\midrule
1 & \xstarfull{3} & 重点掌握基础概念 \\
2 & \xstarfull{4} & 计算题基础入门，重点掌握方法 \\
3 & \xstarfull{4} & 公式多，重点记忆 \\
4 & \xstarfull{5} & \textbf{考试计算题重点}（主要考点，必须掌握计算套路） \\
5 & \xstarfull{5} & \textbf{考试计算题重点}（全书最大的计算章节） \\
6 & \xstarfull{3} & 量纲转化、相似理论，掌握基本概念与公式套用 \\
7 & \xstarfull{3} & 以概念为主 \\
8 & \xstarfull{2} & 选择、填空题考点，掌握基本概念 \\
\bottomrule
\end{tabularx}
\end{center}

\noindent 计算题占总分约 50\%；第 1--3 章偏概念，掌握基本公式与解释即可；
第四章、第五章是计算题主战场；第六至第八章只有少量填空/判断/选择。

\section*{五、重要度星级判定依据（本书严格执行，非主观印象）}

本书对每一个考点、每一道题标注星级，星级只由下列\textbf{文件证据}决定，不使用主观判断。
每条星级后均以〔依据：$\cdots$〕标注到具体文件与页码/题号。

\begin{center}
\begin{tabularx}{\textwidth}{@{}clX@{}}
\toprule
编号 & 证据文件 & 可提取的依据 \\
\midrule
E1 & \texttt{8-1考情分析.pptx} & slide15 题型分值、slide16 各题型考法、slide25 出题风格、slide33 逐章复习规划 \\
E2 & \texttt{9-1考点分析.pptx} & 计算题占 50\%；逐章命题形式（概念/计算） \\
E3 & \texttt{西华大学 818工程流体力学考研分析.docx} & 官方参考书、初试只考前 8 章、真题源于课后习题 \\
E4 & 历年真题 2014--2022（9 套）+ 2026 回忆版 & \textbf{逐题频次统计}：某考点出现的年份与次数 \\
E5 & 题库（试题库（四）、西华选择题题库精简版、补充题库） & 题目来源、与真题的重合情况 \\
E6 & 赵琴《流体力学与流体机械》第 1--8 章 & 章节归属、公式与适用条件 \\
\bottomrule
\end{tabularx}
\end{center}

\begin{center}
\begin{tabularx}{\textwidth}{@{}clX@{}}
\toprule
星级 & 含义 & 判定条件（满足其一） \\
\midrule
\xstarfull{5} & 必考、必须会做 & E4 中作为计算题/大题出现 $\ge 3$ 次；或 E1/E2 明确写"考试计算题重点" \\
\xstarfull{4} & 高频，务必掌握 & E4 中出现 2--3 次；或 E1 写"重点掌握" \\
\xstarfull{3} & 常考，要会 & E4 中出现 1 次；或 E2 写"掌握基本概念与公式套用" \\
\xstarfull{2} & 可能考，了解 & E4 中未出现，但属考纲章节核心内容且 E5 题库反复出现 \\
\xstarfull{1} & 边缘 & 仅复试范围或补充了解（气体动力学、流体机械等） \\
\bottomrule
\end{tabularx}
\end{center}

\noindent\textbf{难度标注}（\nandu{易}/\nandu{中}/\nandu{难}）依据该考点在真题中要求的步骤数：
一步套公式为\emph{易}；需选断面 + 列方程 + 联立为\emph{中}；需积分、求极值或多方程耦联为\emph{难}。

\section*{六、本讲义的使用方法}

\begin{enumerate}
  \item \textbf{第一轮（对照教材）}：按章读"是什么 $\to$ 公式含义 $\to$ 适用条件"，把每个考点框当作教材的最小复习单元；
  \item \textbf{第二轮（对照真题）}：只看〔依据〕里出现过的年份，回真题册核对原题；
  \item \textbf{第三轮（冲刺）}：只看 \xstarfull{4} 以上考点与"易错点"框；
  \item 公式一律\textbf{不背证明}，但必须能背出\textbf{含义}和\textbf{适用条件}——选择题与填空题恰恰考这两点。
\end{enumerate}
